\documentclass[12pt]{article}
\usepackage[utf8]{inputenc}
\usepackage[T1]{fontenc}
\usepackage[margin=1in]{geometry}
\usepackage{setspace}
\usepackage{xcolor}
\usepackage[colorlinks=true, urlcolor=blue!60!black, linkcolor=blue!60!black]{hyperref}
\setlength{\parindent}{0em}
\setlength{\parskip}{0.6em}
\renewcommand{\baselinestretch}{1.0}
\begin{document}
\begin{center}
{\LARGE \textit{Tennessee Wine \& Spirits Retailers Ass'n v. Thomas}, 588 U.S. 504 (2019)\par}
\vspace{0.5em}
{\large Case Brief --- Two-Page Edition\par}
\end{center}

\section{Memory Jogger}
Tennessee's two-year residency rule for liquor retailers violates the dormant Commerce Clause and is not saved by \S~2 of the Twenty-first Amendment (\href{https://www.supremecourt.gov/opinions/18pdf/18-96_5i36.pdf}{Slip Opinion}).

\section{Detailed Case Facts}
Tennessee required two years of residency for an initial retail liquor license, and full resident ownership for corporations. The State's Attorney General thought the rules unconstitutional and the licensing commission stopped enforcing them. When Total Wine, owned by Maryland residents, and the Ketchums, recent arrivals from Utah, sought licenses, an association of incumbent stores threatened suit.

\section{Procedural History}
The federal district court held the rules unconstitutional under \textit{Granholm v. Heald}, and a divided Sixth Circuit agreed. On review of the two-year rule only, the Supreme Court affirmed in 2019.

\section{Judicial Votes}
7--2. Alito wrote for Roberts, Ginsburg, Breyer, Sotomayor, Kagan, and Kavanaugh; Gorsuch dissented, joined by Thomas.

\section{Holding}
The residency rule discriminates against nonresidents, and \S~2 does not shield alcohol laws whose predominant effect is protectionism rather than evidence-backed health or safety.

\section{Analysis of Opinions}
\textbf{Majority:} Section 2 constitutionalized the Wilson and Webb-Kenyon Acts, which let States regulate imported alcohol evenhandedly but not protect local businesses. \textit{Granholm}'s rule reaches retailers, and residency is not an ``essential feature'' of the three-tier system. Tennessee's rationales rested on speculation: agents, background checks, inspections, and revocation would serve them, and the rule was underinclusive. \textbf{Gorsuch:} the people overrode the dormant Commerce Clause for alcohol in \S~2, residency laws are long-standing, and the ``predominant effect'' test is unworkable.

\section{Examples: Future Applications}
\textbf{Same-Side:} a rule licensing only long-time in-state property owners is invalid. \textbf{Opposite-Side:} a state-run retail monopoly excludes all private sellers equally and stands. \textbf{Fence-Sitter:} barring out-of-state retailers without an in-state store from shipping wine, which has divided the circuits.

\section{Critique}
The hybrid standard is indeterminate, and the Court never explains when higher prices are legitimate temperance rather than protectionism. The split pits \S~2's text against the Constitution's anti-protectionist structure.

\section{Key Quotations}
\begin{quote}
``Section 2 gives the States regulatory authority that they would not otherwise enjoy, but as we pointed out in \textit{Granholm}, `mere speculation' or `unsupported assertions' are insufficient to sustain a law that would otherwise violate the Commerce Clause. \ldots{} Where the predominant effect of a law is protectionism, not the protection of public health or safety, it is not shielded by \S~2.''
\end{quote}
\noindent --- Justice Alito, for the Court, 588 U.S. at 539--540.

\end{document}
